% Part: first-order-logic
% Chapter: beyond
% Section: higher-order-logic

\documentclass[../../../include/open-logic-section]{subfiles}

\begin{document}

\olfileid{fol}{byd}{hol}

\olsection{د لوړې درجې منطق}

له لومړۍ درجې منطق څخه دويمې درجې منطق ته په تېرېدو مو وکړے شول چې
په صوري ژبه کښې د لومړۍ درجې د !!{domain} د څيزونو د سټونو په هکله
خبرې وکړو۔ ولې همدلته ودرېږو؟ د بېلګې په توګه، د درېيمې درجې منطق
بايد موږ ته د څيزونو د سټونو د سټونو، يا ښايي آن د هغو سټونو، په
هکله د خبرو وس راکړي چې هم څيزونه او هم د څيزونو سټونه لري۔ د څلورمې
درجې منطق بيا د دې ډول څيزونو د سټونو په هکله خبرې راکوي۔ لکه چې
اټکل مو کړے وي، دا مفکوره په خپله خوښه هر څومره تکرارولے شو۔

په عملي کار کښې د لوړې درجې منطق ډېر ځله د اړيکو پر ځاے د تابعو په
اصطلاحاتو کښې د !!{formula} په بڼه وړاندې کېږي۔ (له طبيعي يو شان ګڼلو
سره دا توپير بنسټيز نۀ دے۔) چې ځينې بنسټيز ``ډولونه'' $A$، $B$،
$C$،~\dots راکړل شي (چې اوس به ئې ``ټايپونه'' بولو)، د لاندې شرط په
ورکولو نوي ټايپونه جوړولے شو:
\begin{quote}
که $\sigma$ او $\tau$ متناهي ټايپونه وي، $\sigma \to \tau$ هم
متناهي ټايپ دے۔
\end{quote}
ټايپونه داسې نحوي ``نښکې'' وګڼئ چې هغه څيزونه ډلبندي کوي چې په خپل
!!{domain} کښې ئې غواړو؛ $\sigma \to \tau$ هغه څيزونه بيانوي چې د
ټايپ~$\sigma$ څيزونه د ټايپ~$\tau$ څيزونو ته وړونکې تابعې دي۔ د
بېلګې په توګه، ښايي د ``رښتوني'' او ``دروغ'' د رښتياوالي قيمتونو يو
ټايپ~$\Omega$ او د طبيعي عددونو يو ټايپ~$\Nat$ وغواړو۔ په دې حالت
کښې د ټايپ~$\Nat \to \Omega$ څيزونه د يوځاييزو اړيکو يا د~$\Nat$
فرعي سټونو په توګه ګڼلے شئ؛ د ټايپ~$\Nat \to \Nat$ څيزونه له طبيعي
عددونو څخه طبيعي عددونو ته تابعې دي؛ او د ټايپ
$(\Nat \to \Nat) \to \Nat$ څيزونه ``فانکشنلونه'' دي، يعنې د لوړې
درجې داسې تابعې چې تابعې اخلي او عددونه ورکوي۔

لکه د دويمې درجې منطق، د لوړې درجې منطق هم د ډېرو ډولونو يو داسې
منطق ګڼلے شو چې د څيز د هر مطلوب ټايپ دپاره يو ډول لري۔ خو اکثره دا
لا روښانه وي چې د لوړو ټايپونو نحو له بنسټه تعريف کړو۔ د بېلګې په
توګه، د متناهي ټايپونو يو سټ په استقرايي ډول داسې تعريفولے شو:
\begin{enumerate}
\item $\Nat$ يو متناهي ټايپ دے۔
\item که $\sigma$ او $\tau$ متناهي ټايپونه وي، $\sigma
  \to \tau$ هم متناهي ټايپ دے۔
\item که $\sigma$ او $\tau$ متناهي ټايپونه وي، $\sigma \times
  \tau$ هم متناهي ټايپ دے۔
\end{enumerate}
په وجداني ډول، $\Nat$ د طبيعي عددونو ټايپ، $\sigma \to \tau$ له
$\sigma$ څخه~$\tau$ ته د تابعو ټايپ، او $\sigma \times \tau$ د هغو
جوړو ټايپ ښيي چې يو څيز ئې له~$\sigma$ او بل له~$\tau$ څخه وي۔ بيا د
ترمونو يو سټ په استقرايي ډول داسې تعريفولے شو:
\begin{enumerate}
\item د هر ټايپ~$\sigma$ دپاره د $x$، $y$، $z$، \dots په څېر د
  ټايپ~$\sigma$ متغيرونو يوه زېرمه شته۔
\item $\Obj 0$ د ټايپ~$\Nat$ يو ترم دے۔
\item $\Obj S$ (ناستي) د ټايپ~$\Nat \to \Nat$ يو ترم دے۔
\item که $s$ د ټايپ~$\sigma$ ترم او $t$ د ټايپ
  $\Nat \to (\sigma \to \sigma)$ ترم وي، $\Obj{R}_{st}$ د ټايپ
  $\Nat \to \sigma$ يو ترم دے۔
\item که $s$ د ټايپ~$\tau \to \sigma$ ترم او $t$ د ټايپ~$\tau$
  ترم وي، $s(t)$ د ټايپ~$\sigma$ يو ترم دے۔
\item که $s$ د ټايپ~$\sigma$ ترم او $x$ د ټايپ~$\tau$ متغير وي،
  $\lambd[x][s]$ د ټايپ~$\tau \to \sigma$ يو ترم دے۔
\item که $s$ د ټايپ~$\sigma$ ترم او $t$ د ټايپ~$\tau$ ترم وي،
  $\tuple{s, t}$ د ټايپ~$\sigma \times \tau$ يو ترم دے۔
\item که $s$ د ټايپ~$\sigma \times \tau$ يو ترم وي، $p_1(s)$ د
  ټايپ~$\sigma$ يو ترم او $p_2(s)$ د ټايپ~$\tau$ يو ترم دے۔
\end{enumerate}
په وجداني ډول، $\Obj{R}_{st}$ هغه تابع ښيي چې په بازګشتي ډول داسې
تعريف شوې وي:
\begin{align*}
\Obj{R}_{st}(0) & = s \\
\Obj{R}_{st}(x+1) & = t(x, R_{st}(x)),
\end{align*}
$\tuple{s, t}$ هغه جوړه ښيي چې لومړے جز ئې~$s$ او دويم جز ئې~$t$
دے، او $p_1(s)$ او~$p_2(s)$ د~$s$ لومړے او دويم غړي
(``پروجيکشنونه'') ښيي۔ په پاى کښې، $\lambd[x][s]$ هغه تابع~$f$ ښيي
چې داسې تعريف شوې وي:
\[
f(x) = s
\]
د ټايپ~$\tau$ د هر~$x$ دپاره؛ نو توکے~(6) د ادراک يوه بڼه راکوي او
د ترمونو په وسيله د تابعو تعريفول شوني کوي۔ !!^{formula}s د يوه شان
ټايپ د ترمونو تر منځ له !!{identity} ويناوو~$\eq[s][t]$، عادي بياني
تړونکو او د لوړو ټايپونو له کميت ټاکنې څخه جوړېږي۔ بيا د نظام بديهي
اصول هغه بنسټيزې معادلې ګڼلے شو چې پر پاس تعريف شويو ترمونو واکمنې
دي، او ورسره د کميت ټاکونکو او !!{identity} لرونکي منطق عادي قاعدې۔
\textbf{د ژباړې د نسخې سمون (OLFOL-033):} په سرچينه کښې د لامبډا ترم
متغير د پايلې ټايپ ته منسوب شوے و؛ دلته هغه د توکي~(6) له تعريف سره
سم د تابع د ورودي ټايپ ته منسوب شوے دے۔

که د متناهي ټايپونو نظام ته د رښتياوالي د قيمتونو ټايپ~$\Omega$
ورزيات شي، بايد د هغه پر کارونه واکمن بديهي اصول هم شامل شي۔ په حقيقت
کښې، که لږ ځير واوسو، پېچلي !!{formula}s په بشپړ ډول لرې کولے او د
ټايپ~$\Omega$ په ترمونو ئې بدلولے شو!{} بيا د ثبوت نظام هم ورسره سم
بدلېدے شي۔ پايله په اصل کښې د \emph{ټايپونو ساده تيوري} ده چې الونزو
چرچ په ۱۹۳۰مو کلونو کښې وړاندې کړې وه۔

لکه د دويمې درجې منطق، د لوړو ټايپونو د معنٰی پوهنې هم بېلابېلې بڼې
شته چې ښايي کارول ئې وغواړو۔ په بشپړه بڼه کښې د ټايپ
$\sigma \to \tau$ متغيرونه د ټايپ~$\sigma$ له څيزونو څخه د
ټايپ~$\tau$ څيزونو ته د \emph{ټولو} تابعو پر سټ ګرځي۔ لکه چې تمه
کېږي، دا معنٰی پوهنه دومره پياوړې ده چې يو بشپړ او اغېزمن
!!{derivation} نظام نۀ مني۔ خو کمزورې معنٰی پوهنه هم څېړلے شو چې
!!a{structure} پکې د هر ټايپ~$\tau$ دپاره د غړو سټونه~$T_\tau$ او
ورسره د تطبيق، پروجيکشن او داسې نورو مناسب عملونه لري۔ که جزئيات سم
بشپړ شي، د پورته بيان شويو !!{derivation} نظامونو د ډولونو دپاره د
بشپړتيا قضيې ترلاسه کولے شو۔

د لوړو ټايپونو منطق ځکه زړه راښکونکے دے چې يو داسې چوکاټ برابروي چې
د رياضياتو لويه برخه پکې په طبيعي ډول ځايولے شو: له~$\Nat$ څخه په
پيل حقيقي عددونه، مسلسل تابعې او داسې نور تعريفولے شو۔ دا منطق د
شهودي منطق په چاپېريال کښې په ځانګړي ډول زړه راښکونکے هم دے، ځکه
ټايپونه روښانه ``تعميري'' تفسيرونه لري۔ په حقيقت کښې د لوړو ټايپونو
د معنٰی پوهنې تعميري بڼې (چې د کلاسیک منطق پر ځاے پر شهودي منطق
ولاړې وي) جوړولے شو؛ دا بڼې دغه تعميري تفسيرونه ډېر ښه روښانوي او په
ډېرو لارو د خپلو کلاسیکو سيالانو په پرتله په زړه پورې دي۔

\end{document}
